\begin{abstract}
Statute retrieval benchmarks typically evaluate at a single query formality level, in a single language, and against citation-derived relevance judgments that over-represent frequently cited provisions. We construct KUHPerdata, the first Indonesian statute retrieval dataset with two query formality levels (case summaries and layperson questions) sharing identical relevance judgments, and expand the ground truth via LLM-based legal element analysis to address citation sparsity. Across KUHPerdata and three existing datasets (BSARD, STARD, IL-PCSR), we find that (i) vocabulary overlap is the primary predictor of lexical retrieval failure across all four languages, (ii) co-relevant statutes cluster structurally while hard negatives are distant, and (iii) sparse citation-based labels create training imbalance where learned methods perform well on frequently cited articles but poorly on the long tail. Based on these findings, we extend the paragraph-level graph neural network of \citet{ilpcsr} with structural node features encoding act membership and positional rank, and post-training alpha grid search, which we call StructGNN. On four datasets spanning four languages, StructGNN improves 24--94\% in MRR@10 over the strongest baseline, with consistent performance across both formality levels. Debiased evaluation confirms that structural features partially mitigate the training imbalance.
\end{abstract}
